\chapter{Conclusion}

This thesis presented the design, development, and evaluation of a browser-based Virtual Reality laboratory for physics education. The primary objective was to investigate whether modern web technologies could be combined with immersive Virtual Reality and real-time collaboration to create an accessible platform for interactive physics learning without requiring the installation of dedicated software.

To achieve this objective, a complete virtual laboratory environment was developed using Babylon.js, WebXR, and Colyseus. The system supports both desktop and immersive Virtual Reality access through standard web browsers, enabling users to explore the laboratory, interact with virtual equipment, and participate in collaborative learning sessions from different devices. The browser-based architecture improves accessibility while maintaining compatibility across multiple hardware platforms.

Two representative physics experiments were implemented to demonstrate the capabilities of the proposed platform. The Convex Lens experiment allows users to investigate image formation by interactively modifying the positions of the lens and projection screen as well as the focal length, while providing real-time visualization based on the thin-lens equation. The Stern--Gerlach experiment illustrates the quantum mechanical concept of spin quantization through an interactive three-dimensional simulation. Together, these experiments demonstrate that the developed platform can effectively support both classical and modern physics education within an immersive virtual environment.

A significant contribution of this thesis is the integration of collaborative multi-user functionality into a browser-based virtual laboratory. Through the client--server architecture, multiple users can simultaneously access the same laboratory session, observe each other's avatars, and interact with shared experimental objects in real time. This collaborative capability extends the virtual laboratory beyond conventional single-user simulations and creates opportunities for group learning, instructor-guided demonstrations, and remote laboratory activities.

The developed system was evaluated through functional testing and a usability study involving eight volunteer participants. The evaluation confirmed that the implemented system successfully satisfied the functional requirements established during the design phase while receiving consistently positive feedback regarding usability, interaction quality, collaborative functionality, and overall learning experience. Participants particularly appreciated the intuitive user interface, immersive Virtual Reality interaction, and the ability to collaboratively explore physics concepts within a shared virtual environment.

Overall, the results of this thesis demonstrate that browser-based Virtual Reality provides a practical and effective approach for developing collaborative virtual laboratories for physics education. By combining immersive visualization, interactive experimentation, and real-time multi-user collaboration within a web-based architecture, the proposed system offers an accessible and extensible platform that can support both classroom teaching and remote learning. The modular design adopted throughout the implementation also provides a strong foundation for future extensions involving additional experiments, enhanced collaborative features, intelligent tutoring, and broader educational deployment.

In conclusion, this work demonstrates the feasibility of integrating modern web technologies, Virtual Reality, and collaborative networking into a unified educational platform. The proposed virtual laboratory not only serves as a proof of concept for immersive browser-based physics education but also establishes a flexible framework that can be expanded to support a wider range of scientific disciplines and future educational applications.